% Viva preparation notes -- compile: latexmk -pdf viva_notes.tex
\documentclass[10.5pt,a4paper]{article}
\usepackage[a4paper,margin=2cm]{geometry}
\usepackage[T1]{fontenc}\usepackage{lmodern}
\usepackage{amsmath,amssymb,mathtools,booktabs,enumitem,microtype}
\usepackage[hidelinks]{hyperref}
\setlist{itemsep=1pt,topsep=2pt}
\newcommand{\PnL}{\mathrm{PnL}_T}
\title{Viva Notes: Delta Hedging of European Options\\\large MAC F244 project (topic 12) \textemdash{} for every group member}
\date{Prepared 13 November 2026}
\begin{document}\maketitle
\noindent Every member should be able to answer \emph{any} question below, not only questions about their own slides.
Numbers quoted are from \texttt{results/full/tables} (seed 2026, 40{,}000 paths).

\tableofcontents

\section{The story in plain English (memorise this)}
We sell a call option. The buyer pays us a premium now; at maturity we owe $\max(S_T-K,0)$. If the stock rises a lot we lose a lot.
To neutralise that risk we buy $\Delta$ shares, where $\Delta$ is how much the option's value moves per unit move in the stock. As the stock moves, $\Delta$ changes, so we must keep
trading. Black--Scholes says: if you trade \emph{continuously}, pay \emph{no costs} and know the \emph{true volatility}, the premium is exactly enough to pay the
option off in every scenario. We relax each of those three assumptions and measure what goes wrong. Coarse trading makes errors \emph{noisy}, at about $1/\sqrt N$. A wrong volatility makes errors
\emph{biased}, and more trading cannot fix that. Costs make frequent trading \emph{expensive}, so there is a best frequency.

\section{Symbols and units}
\begin{tabular}{@{}lll@{}}\toprule
Symbol & Meaning & Units\\\midrule
$S_t$ & stock price at time $t$ & currency (we use $S_0=100$)\\
$K$ & strike & currency\\
$T,\ \tau=T-t$ & maturity, time remaining & years\\
$r$ & continuously compounded risk-free rate & per year\\
$\mu$ & real-world expected return (drift) & per year\\
$\sigma$ & volatility of log returns & per $\sqrt{\text{year}}$\\
$\sigma_h$ & volatility the hedger \emph{uses} & per $\sqrt{\text{year}}$\\
$N,\ \Delta t=T/N$ & number of rebalances, interval & --, years\\
$\phi_k$ & shares held after rebalancing at $t_k$ & shares\\
$B_k$ & cash account after rebalancing at $t_k$ & currency\\
$c$ & proportional cost per side (0.001 = 10 bp) & fraction of traded value\\
$\Delta,\Gamma,\mathcal V$ & $\partial C/\partial S$, $\partial^2C/\partial S^2$, $\partial C/\partial\sigma$ & --, 1/currency, currency per unit $\sigma$\\\bottomrule
\end{tabular}

\section{Mathematics, step by step}
\subsection{GBM and the two measures}
$dS=\mu S\,dt+\sigma S\,dW$ under $\mathbb P$. Apply It\^o to $\ln S$: $d\ln S=(\mu-\tfrac12\sigma^2)dt+\sigma dW$, so
$S_{t+\Delta t}=S_t\exp\{(\mu-\tfrac12\sigma^2)\Delta t+\sigma\sqrt{\Delta t}Z\}$ \emph{exactly}, for any $\Delta t$ (no Euler error).
Under the risk-neutral measure $\mathbb Q$ (Girsanov, with market price of risk $\lambda=(\mu-r)/\sigma$), $dS=rS\,dt+\sigma S\,dW^{\mathbb Q}$. \textbf{Why both?}
Prices are $\mathbb Q$-expectations of discounted payoffs. Hedging P\&L happens in the real world, so we simulate it under $\mathbb P$.
The validation shows that the $\mathbb Q$-Monte-Carlo price matches BS ($z=-0.35$) and that the same estimator under $\mathbb P$ with $\mu=25\%$ does not ($z=401$).

\subsection{It\^o's lemma, replication and the PDE}
For $V(t,S)$: $dV=\big(V_t+\mu SV_S+\tfrac12\sigma^2S^2V_{SS}\big)dt+\sigma SV_S\,dW$.
Hold the option and short $V_S$ shares: $\Pi=V-V_SS$, so $d\Pi=dV-V_SdS=(V_t+\tfrac12\sigma^2S^2V_{SS})dt$. This is riskless, so no-arbitrage requires $d\Pi=r\Pi\,dt$:
\[V_t+rSV_S+\tfrac12\sigma^2S^2V_{SS}-rV=0,\quad V(T,S)=(S-K)^+.\]
$\mu$ has disappeared, and that is why the price does not depend on expected returns. The solution is the BS formula (via Feynman--Kac, $V=e^{-r\tau}\mathbb E^{\mathbb Q}[(S_T-K)^+]$).

\subsection{Price, delta and maturity}
$d_{1,2}=\frac{\ln(S/K)+(r\pm\frac12\sigma^2)\tau}{\sigma\sqrt\tau}$, $C=SN(d_1)-Ke^{-r\tau}N(d_2)$, $\Delta=N(d_1)$, $\Gamma=\varphi(d_1)/(S\sigma\sqrt\tau)$, $\mathcal V=S\varphi(d_1)\sqrt\tau$.
At $\tau\to0$, $d_1\to\pm\infty$ depending on whether $S\gtrless K$, so $\Delta\to\mathbf 1\{S>K\}$ and $C\to(S-K)^+$. The formula divides by $\sqrt\tau$, so the code handles $\tau=0$
separately and never evaluates delta at $T$. Near maturity ATM, $\Gamma\sim1/\sqrt\tau$ blows up: this is where discrete hedging hurts most.
Put--call parity $C-P=S-Ke^{-r\tau}$ is model-free (a forward contract), so it is an independent check.

\subsection{Self-financing bookkeeping (derive this on the board)}
Self-financing means that, between dates, no money enters or leaves except interest. At a rebalance the change in stock holding is paid for in cash:
\begin{align*}
t_0:&\quad B_0=V_0-\phi_0S_0-c|\phi_0|S_0\\
t_k:&\quad B_k=B_{k-1}e^{r\Delta t}-(\phi_k-\phi_{k-1})S_k-c|\phi_k-\phi_{k-1}|S_k\\
T:&\quad \PnL=B_{N-1}e^{r\Delta t}+\phi_{N-1}S_T(1-c)-(S_T-K)^+.
\end{align*}
\textbf{Gains form} ($c=0$): unroll the recursion. $B_{N-1}e^{r\Delta t}=V_0e^{rT}-\sum_k(\phi_k-\phi_{k-1})S_ke^{r(T-t_k)}$; summation by parts gives
$\PnL=V_0e^{rT}+\sum_{k=0}^{N-1}\phi_k\big(S_{k+1}-S_ke^{r\Delta t}\big)e^{r(T-t_{k+1})}-(S_T-K)^+$.
Interpretation: premium grown at $r$, plus the excess return of the stock over cash on each period's holding, minus the payoff. The code computes both forms and they agree to $10^{-13}$.

\textbf{Costs are linear.} Delta ignores $c$, so the trades are identical for all $c$ and $\PnL(c)=\PnL(0)-c\,\mathrm{TV}$, where $\mathrm{TV}$ is the forward value of all traded notional.

\subsection{Worked example (in \texttt{results/validation/manual\_ledger\_cost0.0.csv})}
$K=100$, $T=0.25$, $r=5\%$, $\sigma=20\%$, $N=4$, path $100\to103\to98\to101\to104$. Premium $4.614997$.
\begin{center}\small\begin{tabular}{rrrrrrrr}\toprule
$k$ & $t$ & $S$ & $\Delta$ held & trade & cash for trade & interest & cash after\\\midrule
0&0&100&0.569460&+0.569460&$-56.946018$&0&$-52.331021$\\
1&0.0625&103&0.688948&+0.119488&$-12.307226$&$-0.163790$&$-64.802037$\\
2&0.125&98&0.435666&$-0.253281$&+24.821577&$-0.202823$&$-40.183284$\\
3&0.1875&101&0.612755&+0.177089&$-17.885941$&$-0.125769$&$-58.194993$\\
4&0.25&104&0&$-0.612755$&+63.726513&$-0.182144$&$1.349376$\\\bottomrule
\end{tabular}\end{center}
At $k=4$ we also pay the payoff $4$: $-58.194993\,e^{0.05/16}+0.612755\times104-4=1.349376$. Check one line by hand:
$-52.331021\times(e^{0.003125}-1)=-0.163790$. With $c=10$\,bp the P\&L is $1.172645$; the difference $0.176731$ is $0.001\times\mathrm{TV}$.
\emph{Why positive?} On this path the stock moved less than $\sigma=20\%$ implies on several steps, so the short-gamma position gained. Another path could lose.

\subsection{Why discrete hedging is not exact}
Over $[t_k,t_{k+1}]$, Taylor-expand $C$: $\delta C\approx C_t\Delta t+\Delta\,\delta S+\tfrac12\Gamma(\delta S)^2$. The seller is short the option and long $\Delta$ shares financed at $r$.
Using the PDE to eliminate $C_t$, the one-step P\&L is
\[\epsilon_k\approx\tfrac12\Gamma_kS_k^2\Big[\sigma^2\Delta t-\big(\tfrac{\delta S_k}{S_k}\big)^2\Big]=\tfrac12\Gamma_kS_k^2\sigma^2\Delta t\,(1-Z_k^2).\]
$\mathbb E[1-Z^2]=0$, so there is no bias, and $\operatorname{Var}(1-Z^2)=2$, so $\operatorname{Var}\epsilon_k=\tfrac12\Gamma_k^2S_k^4\sigma^4\Delta t^2$. The steps are uncorrelated, so the variances add:
$\operatorname{Var}\PnL\approx\tfrac12\sigma^4\Delta t\,\mathbb E\!\int_0^T\Gamma^2S^4dt=O(1/N)$. In continuous time, $\sum(\delta S/S)^2\to\sigma^2T$ (quadratic variation), and every $\epsilon_k$ term vanishes.

\textbf{Kamal--Derman.} For ATM, freeze $S\approx S_0$ and use $\Gamma_t\approx\Gamma_0\sqrt{T/(T-t)}$. Then $\sum_k\tfrac12(\Gamma_kS_k^2\sigma^2\Delta t)^2\approx\tfrac{\pi}{4}N(\Gamma_0S_0^2\sigma^2\Delta t)^2$, using $\int_0^T\frac{dt}{\sqrt{T^2-t^2}}=\pi/2$.
With $\mathcal V=S_0^2\Gamma_0\sigma T$ this gives $\mathrm{sd}\approx\sqrt{\pi/4}\,\sigma\mathcal V/\sqrt N$. Base case: $\mathcal V=19.644$, so at $N=63$ the prediction is 0.4387 and the simulation gives 0.4341.

\subsection{Wrong volatility}
Price and hedge at $\sigma_h$. The same step, with $C_t$ eliminated using the PDE at $\sigma_h$, gives $\epsilon_k\approx\tfrac12\Gamma^h_kS_k^2[\sigma_h^2\Delta t-(\delta S/S)^2]$. Its mean is now
$\tfrac12\Gamma^hS^2(\sigma_h^2-\sigma^2)\Delta t\neq0$. Summing and growing to $T$:
$\PnL\approx\tfrac12\int_0^Te^{r(T-t)}\Gamma_h S^2(\sigma_h^2-\sigma^2)\,dt$. The sign depends only on $\sigma_h^2-\sigma^2$ because $\Gamma>0$ for a call. The \emph{size} is path-dependent
(how long the path spends near the strike). That is why the sd does not go to 0 when $N\to\infty$.

\subsection{Transaction costs and Leland}
The expected traded shares per step are $\approx\Gamma S\cdot\mathbb E|\delta S|=\Gamma S^2\sigma\sqrt{\Delta t}\sqrt{2/\pi}$, so the cost per step is $\approx c\,\Gamma S^2\sigma\sqrt{2/\pi}\sqrt{\Delta t}$ and the total cost is $\propto c\sqrt N$.
Leland writes this as an extra $\tfrac12\Gamma S^2\cdot\sigma^2\sqrt{2/\pi}\,k/(\sigma\sqrt{\Delta t})\,\Delta t$ with round-trip $k=2c$, which is the same as using
$\sigma_L^2=\sigma^2(1+\sqrt{2/\pi}\,k/(\sigma\sqrt{\Delta t}))$. Our simulation shows the remaining mean loss under Leland ($0.12$ at 10\,bp) equals the initial-purchase plus liquidation cost ($0.121$). The model does not cover those trades.

\section{Code map (equation $\rightarrow$ function)}
\begin{tabular}{@{}ll@{}}\toprule
Exact GBM step & \texttt{gbm.simulate\_gbm}\\
$C,P,\Delta,\Gamma,\mathcal V$, $\tau=0$ cases & \texttt{black\_scholes.py}\\
Bookkeeping $B_0,B_k,\PnL$; TV; costs & \texttt{hedge.hedge\_short\_call}\\
Gains form (independent check) & \texttt{hedge.pnl\_gains\_form}\\
Hand ledger & \texttt{hedge.hedge\_ledger}\\
Gamma formula for wrong $\sigma$ & \texttt{hedge.predicted\_pnl\_misspecified}\\
Leland $\sigma_L$ & \texttt{black\_scholes.leland\_sigma}\\
Statistics and SEs & \texttt{stats.summarize}\\
Experiments A--D / E & \texttt{experiments.py} / \texttt{realdata.py}\\\bottomrule
\end{tabular}

\noindent\textbf{Design choices you may be asked to justify.} \emph{Fine grid + subsampling}: every frequency uses the same Brownian path, so comparisons are paired (common random numbers).
\emph{Batches of 5{,}000}: limit memory, and they also serve as 8 independent replications (slope SE). \emph{SeedSequence.spawn}: independent streams per experiment from one master seed.
\emph{Delta-method SE for the sd}: $\mathrm{SE}(s)\approx\sqrt{(m_4-s^4)/n}/(2s)$ uses the sample fourth moment, so it remains valid for fat tails. \emph{Bootstrap} for quantiles and ES.

\section{Key results to quote}
\begin{itemize}
\item A: sd 1.856 ($N=3$) $\to$ 0.434 ($N=63$, daily) $\to$ 0.098 ($N=1260$). Slope $-0.488\pm0.001$. Kamal--Derman within 1.4\% for $N\ge63$. Daily sd is 9.4\% of the premium 4.615.
\item A, drift: sd 0.439/0.434/0.418 for $\mu=5/10/25\%$ at $N=63$. The effect is small but detectable (discrete-time only).
\item B: means $-1.94,-0.98,0,+0.99,+1.98$ for $\sigma_h=10..30\%$ (theory identical to 2 dp). The bias does not shrink with $N$. If the premium is charged at the true $\sigma$, the mean is about 0.
\item C: RMSE-optimal $N$ = 1260/252/126/63/21 for 0/5/10/20/50\,bp. Leland mean $-0.12$ at 10\,bp for all $N$, equal to the setup and unwind cost.
\item D: ATM sd 0.32 vs 0.98 (15\% vs 45\% vol). OTM ($K=120$) sd/premium 254\% vs 30\%. Absolute sd flat in maturity.
\item E: IBM 29\%, MSFT 34\%, GOOG 39\%, AAPL 55\%, AMZN 59\% (2000--10 monthly); AAPL 24\% (2015--17 daily). Backtest corr(P\&L, variance gap) = 0.87.
\end{itemize}

\section{Question bank with model answers}
\begin{enumerate}[label=\textbf{Q\arabic*.},leftmargin=2.2em]
\item \textbf{Why does the option price not depend on $\mu$?} The hedged portfolio is riskless, so it must earn $r$. $\mu$ cancels in $d\Pi$. Risk preferences are absorbed by replication.
\item \textbf{Then why simulate under $\mathbb P$?} The P\&L is realised along real-world paths. In continuous time the result is $\mu$-independent. In discrete time $\mu$ has a small effect (sd falls 4.6\% from $\mu=5\%$ to $25\%$) because paths drift away from the strike, where gamma is lower.
\item \textbf{Is the exact GBM step an approximation?} No. $\ln S$ increments are exactly normal, so any step size is exact in distribution. Validation: $E[S_T]$ and $\operatorname{Var}\ln S_T$ match theory with 1 step and with 252 steps.
\item \textbf{What does self-financing mean in your code?} The only cash flows are the premium at 0, trades, costs, interest and the final payoff. Line \texttt{cash -= trade * S\_k} is exactly that.
\item \textbf{Why is the P\&L at $T$ and not at 0?} So that all cash flows are in the same time-$T$ money. Discounting by $e^{-rT}$ (a factor of 0.988) changes nothing qualitatively.
\item \textbf{Sign convention?} Positive means the seller gains.
\item \textbf{Why does the mean error vanish but the sd does not?} Each step's error is $\propto(1-Z^2)$, which has mean 0 and variance 2.
\item \textbf{Why $1/\sqrt N$?} There are $N$ independent errors, each of size $\propto\Delta t=T/N$, so the variance is $N\cdot(T/N)^2\propto1/N$.
\item \textbf{Your slope is $-0.488$, not $-0.5$. Is your code wrong?} No. $-\tfrac12$ is the asymptotic slope. Coarse grids ($N=3,9$) are outside the asymptotic regime: the Kamal--Derman approximation overstates the sd by 8\% at $N=3$. The slope fitted on $N\ge63$ alone is $-0.497$.
\item \textbf{Is every error statistic monotone in $N$?} Not necessarily, and $|$mean$|$ is not in our sample. We report monotonicity as an observation, not a theorem.
\item \textbf{Why are the tails fat?} The error is a sum of $\Gamma$-weighted $(1-Z^2)$ terms, and the weights are random and large near the strike close to expiry. Paths that end near $K$ dominate the losses.
\item \textbf{Why is $N=3$ slightly biased?} The second-order Taylor expansion fails for large moves, and higher-order terms do not average out. The bias is small ($-0.022\pm0.009$).
\item \textbf{Why is over-estimating $\sigma$ good for the seller?} The seller collects a larger premium, and the gamma formula shows a positive expected P\&L whenever $\sigma_h>\sigma$ (El Karoui et al.\ 1998: a super-hedge if the hedge vol dominates).
\item \textbf{If the premium is fair but the delta is wrong?} The mean is about 0, but the sd is unchanged and does not vanish as $N\to\infty$. The wrong delta adds risk, not expected loss.
\item \textbf{Why do costs grow like $\sqrt N$?} The shares traded per step are $\propto\Gamma|\delta S|\propto\sqrt{\Delta t}$, and there are $N$ steps, so the total is $\propto N\sqrt{T/N}=\sqrt{NT}$.
\item \textbf{Why is your Leland mean not zero?} Leland's correction prices only \emph{rebalancing} costs. We also charge the initial purchase and the final liquidation ($0.121$ at 10\,bp), which explains the residual exactly.
\item \textbf{Why $k=2c$?} Leland defines $k$ as the round-trip cost. We charge $c$ per side.
\item \textbf{Is RMSE the right objective?} It is one convention that combines bias and spread. A risk-averse seller might minimise mean + $\lambda\cdot$sd, or ES. The optimum would move.
\item \textbf{How do you know your Monte Carlo is accurate enough?} Every statistic carries an SE. With 40{,}000 paths the sd at $N=63$ is $0.434\pm0.002$. Two independent seeds agree within 3 SE.
\item \textbf{Why do OTM options have huge relative errors?} The premium is tiny (0.03), but if the stock reaches the strike, gamma becomes large. The absolute error is small, but relative to the premium it is large.
\item \textbf{Why is the absolute error flat in maturity at fixed daily rebalancing?} $\mathrm{sd}\approx\sqrt{\pi/4}\,\sigma\mathcal V/\sqrt N$, with $\mathcal V\propto\sqrt T$ (ATM) and $N\propto T$, so $T$ cancels.
\item \textbf{Is your real data `market data'?} It is secondary-source historical prices (vega-datasets and plotly/datasets mirrors). We used no option quotes. The premiums are model prices. The results are illustrations, not trading evidence.
\item \textbf{Why are overlapping windows a problem?} Consecutive windows share 62 of 63 days, so the 380 windows are far from independent. Only 7 are non-overlapping.
\item \textbf{Biggest limitation?} Constant volatility. Real volatility is stochastic and jumps exist, and then delta hedging alone cannot replicate the option.
\item \textbf{What would you do next?} Stochastic volatility (Heston); hedging with the Whalley--Wilmott band instead of fixed-time rebalancing; using real option quotes; adding gamma hedging with a second option.
\end{enumerate}

\section{Honest limits of what the plots show}
Fig.\,1 shows the error distribution only in the BS world. Fig.\,2b shows agreement with a continuous-time approximation, not proof of it. Fig.\,3's optimum depends on choosing RMSE.
Fig.\,4 uses two stylised regimes, not calibrated assets. Fig.\,5 is one stock, two years, model premiums, and overlapping windows.
\end{document}
